\chapter{2009年湖北卷（理）}

\section{单选题}

\begin{problem}
已知 \(P = \{\bs{a} \mid \bs{a} = (1,0) + m(0,1), m \in \R\}\)，\(Q = \{\bs{b} \mid \bs{b} = (1,1) + n(-1,1), n \in \R\}\) 是两个向量集合，则 \(P \cap Q =\)（\quad）

\choices
    {\(\{(1,1)\}\)}
    {\(\{(-1,1)\}\)}
    {\(\{(1,0)\}\)}
    {\(\{(0,1)\}\)}
\end{problem}

\begin{answer}
A
\end{answer}

\begin{problem}
函数 \( y = \frac{1 - ax}{1 + ax} \)（\(x \in \R\)，\(x \neq -\frac{1}{a}\)）的反函数是（\quad）

\choices
    {\( y = \frac{1 - ax}{1 + ax} \)（\(x \in \R\)，\(x \neq -\frac{1}{a}\)）}
    {\( y = \frac{1 + ax}{1 - ax} \)（\(x \in \R\)，\(x \neq \frac{1}{a}\)）}
    {\( y = \frac{1 + x}{a(1 - x)} \)（\(x \in \R\)，\(x \neq 1\)）}
    {\( y = \frac{1 - x}{a(1 + x)} \)（\(x \in \R\)，\(x \neq -1\)）}
\end{problem}

\begin{answer}
D
\end{answer}

\begin{problem}
投掷两颗骰子，得到其向上的点数分别为 \(m\) 和 \(n\)，则复数 \( (m + n\i)(n - m\i) \) 为实数的概率为（\quad）

\choices
    {\(\frac{1}{3}\)}
    {\(\frac{1}{4}\)}
    {\(\frac{1}{6}\)}
    {\(\frac{1}{12}\)}
\end{problem}

\begin{answer}
C
\end{answer}

\begin{problem}
函数 \( y = \cos \left(2x + \frac{\pi}{6}\right) - 2 \) 的图象 \(F\) 按向量 \(\bs{a}\) 平移到 \(F'\)，\(F'\) 的解析式 \( y = f(x) \)，当 \( y = f(x) \) 为奇函数时，向量 \(\bs{a}\) 可以等于（\quad）

\choices
    {\(\left(-\frac{\pi}{6}, -2\right)\)}
    {\(\left(-\frac{\pi}{6}, 2\right)\)}
    {\(\left(\frac{\pi}{6}, -2\right)\)}
    {\(\left(\frac{\pi}{6}, 2\right)\)}
\end{problem}

\begin{answer}
B
\end{answer}

\begin{problem}
将甲，乙，丙，丁四名学生分到三个不同的班，每个班至少分到一名学生，且甲，乙两名学生不能分到一个班，则不同分法的种数为（\quad）

\choices
    {\(18\)}
    {\(24\)}
    {\(30\)}
    {\(36\)}
\end{problem}

\begin{answer}
C
\end{answer}

\begin{problem}
设 \(\left(\frac{\sqrt{2}}{2} + x\right)^{2n} = a_0 + a_1x + a_2x^2 + \cdots + a_{2n-1}x^{2n-1} + a_{2n}x^{2n}\)，则

\(\displaystyle\lim_{n\to \infty}\left[(a_0 + a_2 + a_4 + \cdots + a_{2n})^2 - (a_1 + a_3 + a_5 + \cdots + a_{2n-1})^2\right] = (\quad)\)

\choices
    {\(-1\)}
    {\(0\)}
    {\(1\)}
    {\(\frac{\sqrt{2}}{2}\)}
\end{problem}

\begin{answer}
B
\end{answer}

\begin{problem}
已知双曲线 \(\frac{x^{2}}{2}-\frac{y^{2}}{2}=1\) 的准线过椭圆 \(\frac{x^{2}}{4}+\frac{y^{2}}{b^{2}}=1\) 的焦点，则直线 \(y=kx+2\) 与椭圆至多有一个交点的充要条件是（\quad）

\choices
    {\(k \in \left[-\frac{1}{2}, \frac{1}{2}\right]\)}
    {\(k \in \left(-\infty, -\frac{1}{2}\right] \cup \left[\frac{1}{2}, +\infty\right)\)}
    {\(k \in \left[-\frac{\sqrt{2}}{2}, \frac{\sqrt{2}}{2}\right]\)}
    {\(k \in \left(-\infty, -\frac{\sqrt{2}}{2}\right] \cup \left[\frac{\sqrt{2}}{2}, +\infty\right)\)}
\end{problem}

\begin{answer}
A
\end{answer}

\begin{problem}
在“家电下乡”活动中，某厂要将 \(100\text{ 台}\)洗衣机运往邻近的乡镇，现有 \(4\text{ 辆}\)甲型货车和 \(8\text{ 辆}\)乙型货车可供使用，每辆甲型货车运输费用 \(400\text{ 元}\)，可装洗衣机 \(20\text{ 台}\)；每辆乙型货车运输费用 \(300\text{ 元}\)，可装洗衣机 \(10\text{ 台}\)，若每辆至多只运一次，则该厂所花的最少运输费用为（\quad）

\choices
    {\(2000\text{ 元}\)}
    {\(2200\text{ 元}\)}
    {\(2400\text{ 元}\)}
    {\(2800\text{ 元}\)}
\end{problem}

\begin{answer}
B
\end{answer}

\begin{problem}
设球的半径为时间 \(t\) 的函数 \(R(t)\). 若球的体积以均匀速度 \(c\) 增长，则球的表面积的增长速度与球半径（\quad）

\choices
    {成正比，比例系数为 \(c\)}
    {成正比，比例系数为 \(2c\)}
    {成反比，比例系数为 \(c\)}
    {成反比，比例系数为 \(2c\)}
\end{problem}

\begin{answer}
D
\end{answer}

\begin{problem}
古希腊人常用小石子在沙滩上摆成各种形状来研究数，例如：

他们研究过图1中的 \(1\)，\(3\)，\(6\)，\(10\)，\(\cdots\)，由于这些数能够表示成三角形，将其称为三角形数；类似地，图2中的 \(1\)，\(4\)，\(9\)，\(16\cdots\) 这样的数称为正方形数. 下列数中既是三角形数又是正方形数的是（\quad）

\problemresetlayout
\examfiguregroup[float=inline]{img/2009/hubei_science/q10_fig1.png}{%
    \makebox[0.42\linewidth][c]{\bitmapinclude[max width=\linewidth]{img/2009/hubei_science/q10_fig1.png}}\hspace{1.4em}%
    \makebox[0.42\linewidth][c]{\bitmapinclude[max width=\linewidth]{img/2009/hubei_science/q10_fig2.png}}%
}

\choices
    {\(289\)}
    {\(1024\)}
    {\(1225\)}
    {\(1378\)}
\end{problem}

\begin{answer}
C
\end{answer}

\section{填空题}

\begin{problem}
已知关于 \(x\) 的不等式 \(\frac{ax - 1}{x + 1} < 0\) 的解集是 \((-\infty, -1) \cup \left(-\frac{1}{2}, +\infty\right)\)，则 \(a = \)\fillinblank{}.
\end{problem}

\begin{answer}
\(-2\)
\end{answer}

\begin{problem}
\noindent 样本容量为 \(200\) 的频率分布直方图.
根据样本的频率分布直方图估计，样本数据在 \([6, 10)\) 内的频数为\fillinblank{}，数据落在 \([2, 10)\) 内的概率约为\fillinblank{}.

\bitmapinlinefigure[0.34\linewidth][max width=0.34\linewidth]{img/2009/hubei_science/q12_fig1.png}
\end{problem}

\begin{answer}
64，0.4
\end{answer}

\begin{problem}
卫星和地面之间的电视信号沿直线传输，电视信号能够传送到达的地面区域，称为这个卫星的覆盖区域. 为了转播 2008 年北京奥运会，我国发射了“中星九号”广播电视直播卫星，它离地球表面的距离约为 \(36000\mathrm{~km}\). 已知地球半径约为 \(6400\mathrm{~km}\)，则“中星九号”覆盖区域内的任意两点的球面距离的最大值约为\fillinblank{}\(\mathrm{km}\).（结果中保留反余弦的符号）
\end{problem}

\begin{answer}
\(12800\arccos\frac{8}{53}\)
\end{answer}

\begin{problem}
已知函数 \(f(x) = f'\left(\frac{\pi}{4}\right)\cos x + \sin x\)，则 \(f\left(\frac{\pi}{4}\right)\) 的值为\fillinblank{}.
\end{problem}

\begin{answer}
\(1\).
\end{answer}

\begin{problem}
已知数列 \(\{a_{n}\}\) 满足：\(a_1 = m\)（\(m\) 为正整数），\(a_{n+1} = \begin{cases}
\frac{a_n}{2}, & \text{当 } a_n \text{ 为偶数时},\\
3a_n + 1, & \text{当 } a_n \text{ 为奇数时}.
\end{cases}\) 若 \(a_6 = 1\)，则 \(m\) 所有可能的取值为\fillinblank{}.
\end{problem}

\begin{answer}
\(4,5,32\)
\end{answer}

\section{解答题}

\begin{problem}
一个盒子里装有 \(4\text{ 张}\)大小形状完全相同的卡片，分别标有数 \(2\)，\(3\)，\(4\)，\(5\)；另一个盒子里也装有 \(4\text{ 张}\)大小形状完全相同的卡片，分别标有数 \(3\)，\(4\)，\(5\)，\(6\). 现从一个盒子中任取一张卡片，其上面的数记为 \(x\)；再从另一盒子里任取一张卡片，其上面的数记为 \(y\)，记随机变量 \(\eta = x + y\)，求 \(\eta\) 的分布列和数学期望.
\end{problem}

\begin{solution}
从两个盒子中各取一张卡片时，所有有序取法共有 \(4\times4=16\) 种，且每种取法等可能. 按 \(x+y\) 的值统计取法数，得到如下分布列.

\begin{center}
\begin{tabular}{c|ccccccc}
\hline
\(\eta\) & 5 & 6 & 7 & 8 & 9 & 10 & 11 \\
\hline
取法数 & 1 & 2 & 3 & 4 & 3 & 2 & 1 \\
\hline
\(P(\eta)\) & \(\frac{1}{16}\) & \(\frac{2}{16}\) & \(\frac{3}{16}\) & \(\frac{4}{16}\) & \(\frac{3}{16}\) & \(\frac{2}{16}\) & \(\frac{1}{16}\) \\
\hline
\end{tabular}
\end{center}

因此
\[
E(\eta)=\frac{5+2\times6+3\times7+4\times8+3\times9+2\times10+11}{16}=8.
\]
\end{solution}

\begin{problem}
已知向量 \(\bs{a} = (\cos \alpha, \sin \alpha)\)，\(\bs{b} = (\cos \beta, \sin \beta)\)，\(\bs{c} = (-1, 0)\).

\begin{enumerate}
\item 求向量 \(\bs{b} + \bs{c}\) 的长度的最大值；
\item 设 \(\alpha = \frac{\pi}{4}\)，且 \(\bs{a} \perp (\bs{b} + \bs{c})\)，求 \(\cos \beta\) 的值.
\end{enumerate}
\end{problem}

\begin{solution}
\begin{enumerate}
\item
\[
\left\lvert\bs{b}+\bs{c}\right\rvert^{2}=(\cos\beta-1)^{2}+\sin^{2}\beta=2-2\cos\beta.
\]
因为 \(\cos\beta\geqslant-1\)，所以
\[
\left\lvert\bs{b}+\bs{c}\right\rvert^{2}\leqslant4,
\]
从而向量 \(\bs{b}+\bs{c}\) 的长度的最大值为 \(2\)，且在 \(\cos\beta=-1\) 时取得.

\item 当 \(\alpha=\frac{\pi}{4}\) 时，\(\bs{a}=\left(\frac{\sqrt{2}}{2},\frac{\sqrt{2}}{2}\right)\). 由垂直条件，得
\[
0=\bs{a}\cdot(\bs{b}+\bs{c})=\frac{\sqrt{2}}{2}(\cos\beta-1+\sin\beta),
\]
即 \(\cos\beta+\sin\beta=1\). 又因为
\[
(\cos\beta+\sin\beta)^{2}=\cos^{2}\beta+\sin^{2}\beta+2\sin\beta\cos\beta,
\]
所以 \(\sin\beta\cos\beta=0\). 结合 \(\cos\beta+\sin\beta=1\)，得到 \((\cos\beta,\sin\beta)=(1,0)\) 或 \((\cos\beta,\sin\beta)=(0,1)\).
故 \(\cos\beta=1\) 或 \(0\).
\end{enumerate}
\end{solution}

\begin{problem}
如图，四棱锥 \(S-ABCD\) 的底面是正方形，\(SD \perp\) 平面 \(ABCD\)，\(SD = 2a\)，\(AD = \sqrt{2}a\)，点 \(E\) 是 \(SD\) 上的点，且 \(DE = \lambda a\)（\(0 < \lambda \leqslant 2\)）.

\begin{enumerate}
\item 求证：对任意的 \(\lambda \in (0,2]\)，都有 \(AC \perp BE\)；
\item 设二面角 \(C-AE-D\) 的大小为 \(\theta\). 直线 \(BE\) 与平面 \(ABCD\) 所成的角为 \(\varphi\). 若 \(\tan \theta \cdot \tan \varphi = 1\)，求 \(\lambda\) 的值.
\end{enumerate}

\bitmapfigure[max width=0.34\linewidth]{img/2009/hubei_science/q18_fig1.png}
\end{problem}

\begin{solution}
设正方形边长为 \(s=\sqrt{2}a\)，取 \(D\) 为原点，\(DA\)，\(DC\)，\(DS\) 分别为 \(x\)，\(y\)，\(z\) 轴正向，并设 \(h=\lambda a\). 则 \(A=(s,0,0)\)，\(B=(s,s,0)\)，\(C=(0,s,0)\)，\(D=(0,0,0)\)，\(E=(0,0,h)\).

\begin{enumerate}
\item
有 \(\overrightarrow{AC}=(-s,s,0)\)，\(\overrightarrow{BE}=(-s,-s,h)\).
从而
\[
\overrightarrow{AC}\cdot\overrightarrow{BE}=s^{2}-s^{2}=0.
\]
因此对任意 \(\lambda\in(0,2]\)，都有 \(AC\perp BE\).

\item 平面 \(AEC\) 与平面 \(AED\) 的法向量可分别取为 \(\bs{n}_{1}=\overrightarrow{AE}\times\overrightarrow{AC}=(-hs,-hs,-s^{2})\)，\(\bs{n}_{2}=\overrightarrow{AD}\times\overrightarrow{AE}=(0,sh,0)\).
设二面角 \(C-AE-D\) 的大小为 \(\theta\)，则
\[
\cos\theta=\frac{|\bs{n}_{1}\cdot\bs{n}_{2}|}{|\bs{n}_{1}|\,|\bs{n}_{2}|}
=\frac{h}{\sqrt{2h^{2}+s^{2}}}.
\]
于是
\[
\tan\theta=\frac{\sqrt{h^{2}+s^{2}}}{h}=\frac{\sqrt{\lambda^{2}+2}}{\lambda}.
\]

直线 \(BE\) 在底面 \(ABCD\) 上的射影为 \(BD\)，且 \(BD=\sqrt{2}s=2a\)，所以
\[
\tan\varphi=\frac{DE}{BD}=\frac{\lambda a}{2a}=\frac{\lambda}{2}.
\]
由题设
\[
1=\tan\theta\tan\varphi=\frac{\sqrt{\lambda^{2}+2}}{2}.
\]
因 \(\lambda>0\)，解得 \(\lambda^{2}=2\)，故 \(\lambda=\sqrt{2}\).
\end{enumerate}
\end{solution}

\begin{problem}
已知 \(\{a_{n}\}\) 的前 \(n\) 项和 \(S_{n} = -a_{n} - \left(\frac{1}{2}\right)^{n-1} + 2\)（\(n\) 为正整数）.

\begin{enumerate}
\item 令 \(b_{n} = 2^{n}a_{n}\)，求证数列 \(\{b_{n}\}\) 是等差数列，并求数列 \(\{a_{n}\}\) 的通项公式；
\item 令 \(c_{n} = \frac{n + 1}{n} a_{n}\)，\(T_{n} = c_{1} + c_{2} + \cdots + c_{n}\)，试比较 \(T_{n}\) 与 \(\frac{5n}{2n + 1}\) 的大小，并予以证明.
\end{enumerate}
\end{problem}

\begin{solution}
\begin{enumerate}
\item 当 \(n=1\) 时，由已知式得
\[
a_{1}=S_{1}=-a_{1}-1+2,
\]
所以 \(a_{1}=\frac12\). 当 \(n\geqslant2\) 时，由 \(S_{n}-S_{n-1}=a_{n}\)，得
\[
a_{n}=-a_{n}+a_{n-1}+\left(\frac12\right)^{n-1},
\]
即
\[
2a_{n}=a_{n-1}+\left(\frac12\right)^{n-1}.
\]
令 \(b_{n}=2^{n}a_{n}\)，上式两边乘以 \(2^{n}\)，得
\[
2b_{n}=2b_{n-1}+2,
\]
即 \(b_{n}-b_{n-1}=1\). 又 \(b_{1}=2a_{1}=1\)，所以 \(b_{n}=n\)，从而
\[
a_{n}=\frac{n}{2^{n}}.
\]

\item 由定义
\[
c_{n}=\frac{n+1}{n}a_{n}=\frac{n+1}{2^{n}},
\]
故
\[
T_{n}=\sum_{k=1}^{n}\frac{k+1}{2^{k}}.
\]
记 \(U_{n}=\sum_{k=1}^{n}\frac{k}{2^{k}}\). 将 \(U_n\) 乘以 \(2\) 并错位相减，得
\[
2U_{n}-U_{n}=1+\sum_{k=1}^{n-1}\frac{1}{2^{k}}-\frac{n}{2^{n}}
=2-\frac{n+2}{2^{n}},
\]
所以 \(U_{n}=2-\frac{n+2}{2^{n}}\). 又
\[
\sum_{k=1}^{n}\frac{1}{2^{k}}=1-\frac{1}{2^{n}},
\]
从而
\[
T_{n}=3-\frac{n+3}{2^{n}}.
\]
比较两者，得
\[
T_{n}-\frac{5n}{2n+1}
=(n+3)\left(\frac{1}{2n+1}-\frac{1}{2^{n}}\right)
=\frac{(n+3)\left(2^{n}-2n-1\right)}{(2n+1)2^{n}}.
\]
当 \(n=1,2\) 时，\(2^{n}-2n-1<0\)，故 \(T_{n}<\frac{5n}{2n+1}\). 当 \(n=3\) 时，\(2^{3}>2\times3+1\)；若 \(2^{n}>2n+1\)，则
\[
2^{n+1}>4n+2>2n+3=2(n+1)+1.
\]
因此当 \(n\geqslant3\) 时，\(2^{n}>2n+1\)，故 \(T_{n}>\frac{5n}{2n+1}\).
\end{enumerate}
\end{solution}

\begin{problem}
过抛物线 \(y^{2} = 2px\)（\(p > 0\)）的对称轴上一点 \(A(a,0)\)（\(a > 0\)）的直线与抛物线相交于 \(M\)，\(N\) 两点，自 \(M\)，\(N\) 向直线 \(l: x = -a\) 作垂线，垂足分别为 \(M_{1}\)，\(N_{1}\).

\begin{enumerate}
\item 当 \(a = \frac{p}{2}\) 时，求证： \(AM_{1} \perp AN_{1}\)；
\item 记 \(\triangle AMM_{1}\)，\(\triangle AM_{1}N_{1}\)，\(\triangle ANN_{1}\) 的面积分别为 \(S_{1}\)，\(S_{2}\)，\(S_{3}\)，是否存在 \(\lambda\)，使得对任意的 \(a > 0\)，都有 \(S_{2}^{2} = \lambda S_{1}S_{3}\) 成立. 若存在，求出 \(\lambda\) 的值；若不存在，说明理由.
\end{enumerate}
\end{problem}

\begin{solution}
设直线 \(MN\) 不是竖直线，记其非零斜率为 \(k\)，并设 \(M=\left(\frac{u^{2}}{2p},u\right)\)，\(N=\left(\frac{v^{2}}{2p},v\right)\).
由 \(M\)，\(N\)，\(A(a,0)\) 共线，直线可写为 \(y=k(x-a)\). 代入抛物线方程，得
\[
ku^{2}-2pu-2pak=0,
\]
因此由韦达定理得 \(u+v=\frac{2p}{k}\)，\(uv=-2pa\).

\begin{enumerate}
\item 由 \(M_{1}=(-a,u)\)，\(N_{1}=(-a,v)\)，有 \(\overrightarrow{AM_{1}}=(-2a,u)\)，\(\overrightarrow{AN_{1}}=(-2a,v)\).
所以
\[
\overrightarrow{AM_{1}}\cdot\overrightarrow{AN_{1}}=4a^{2}+uv
=4a^{2}-2pa=2a(2a-p).
\]
当 \(a=\frac{p}{2}\) 时，上式等于 \(0\)，故 \(AM_{1}\perp AN_{1}\).

\item 由 \(uv=-2pa<0\)，可知 \(u\)，\(v\) 异号. 三角形面积分别为 \(S_{1}=\frac12\left(\frac{u^{2}}{2p}+a\right)|u|\)，\(S_{3}=\frac12\left(\frac{v^{2}}{2p}+a\right)|v|\)，以及 \(S_{2}=\frac12\cdot2a\cdot|u-v|=a|u-v|\).
又因为 \(2pa=-uv\)，所以
\[
\left(\frac{u^{2}}{2p}+a\right)\left(\frac{v^{2}}{2p}+a\right)
=\frac{u(u-v)}{2p}\cdot\frac{v(v-u)}{2p}
=\frac{a(u-v)^{2}}{2p}.
\]
于是
\[
4S_{1}S_{3}
=\left(\frac{u^{2}}{2p}+a\right)\left(\frac{v^{2}}{2p}+a\right)|uv|
=\frac{a(u-v)^{2}}{2p}\cdot2pa
=a^{2}(u-v)^{2}=S_{2}^{2}.
\]

若直线 \(MN\) 为竖直线，则 \(x=a\)，且 \(u^{2}=v^{2}=2pa\)，\(v=-u\). 此时 \(S_{1}=a|u|\)，\(S_{3}=a|v|\)，\(S_{2}=a|u-v|=2a|u|\).
仍有 \(S_{2}^{2}=4S_{1}S_{3}\). 因此对任意 \(a>0\) 都成立，存在的常数为 \(\lambda=4\).
\end{enumerate}
\end{solution}

\begin{problem}
在 \(\R\) 上定义运算 \(\otimes\)： \(p \otimes q = -\frac{1}{3}(p - c)(q - b) + 4bc\)（\(p\)，\(q\) 为实常数）. 记 \(f_1(x) = x^{2} - 2c\)，\(f_2(x) = x - 2b\)，\(x \in \R\). 令 \(f(x) = f_1(x) \otimes f_2(x)\).

\begin{enumerate}
\item 如果函数 \(f(x)\) 在 \(x = 1\) 处有极值 \(-\frac{4}{3}\)，试确定 \(b\)，\(c\) 的值；
\item 求曲线 \(y = f(x)\) 上斜率为 \(c\) 的切线与该曲线的公共点；
\item 记 \(g(x) = |f'(x)|\)（\(-1 \leqslant x \leqslant 1\)）的最大值为 \(M\). 若 \(M \geqslant k\) 对任意的 \(b\)，\(c\) 恒成立，试求 \(k\) 的最大值.
\end{enumerate}
\end{problem}

\begin{solution}
由定义
\[
f(x)=-\frac13\bigl((x^{2}-2c)-c\bigr)\bigl((x-2b)-b\bigr)+4bc
=-\frac13(x^{2}-3c)(x-3b)+4bc
=-\frac13x^{3}+bx^{2}+cx+bc.
\]

\begin{enumerate}
\item 由
\[
f'(x)=-x^{2}+2bx+c,
\]
函数在 \(x=1\) 处有极值，必有 \(f'(1)=0\)，即
\[
2b+c=1.
\]
又由 \(f(1)=-\frac43\)，得
\[
-\frac13+b+c+bc=-\frac43,
\]
即 \(b+c+bc=-1\). 代入 \(c=1-2b\)，得
\[
1-2b^{2}=-1,
\]
所以 \(b=1\) 或 \(b=-1\)，相应地 \(c=-1\) 或 \(c=3\).

当 \((b,c)=(1,-1)\) 时，\(f'(x)=-(x-1)^{2}\)，在 \(x=1\) 的两侧导数均不变号，故 \(x=1\) 不是极值点. 当 \((b,c)=(-1,3)\) 时，
\[
f'(x)=-(x-1)(x+3),
\]
导数在 \(x=1\) 的左侧为正、右侧为负，故 \(x=1\) 为极大值点. 因此 \(b=-1\)，\(c=3\).

\item 在第（1）问所得 \(b=-1\)，\(c=3\) 的条件下，\(f(x)=-\frac13x^{3}-x^{2}+3x-3\)，\(f'(x)=-x^{2}-2x+3\). 设切点横坐标为 \(t\). 切线斜率为 \(c=3\)，所以 \(f'(t)=3\)，即 \(-t^{2}-2t+3=3\).
即
\[
t(t+2)=0.
\]
因此切点横坐标为 \(t=0\) 或 \(t=-2\).

当 \(t=0\) 时，\(f(0)=-3\)，切线为
\[
y=3x-3.
\]
将其与曲线方程联立，得
\[
-\frac13x^{3}-x^{2}+3x-3=3x-3,
\]
即
\[
-\frac13x^{2}(x+3)=0.
\]
所以公共点为 \((0,-3)\) 和 \((-3,-12)\).

当 \(t=-2\) 时，\(f(-2)=-\frac{31}{3}\)，切线为
\[
y=3(x+2)-\frac{31}{3}=3x-\frac{13}{3}.
\]
将其与曲线方程联立，得
\[
-\frac13x^{3}-x^{2}+3x-3=3x-\frac{13}{3},
\]
即
\[
-\frac13(x+2)^{2}(x-1)=0.
\]
所以公共点为 \(\left(-2,-\frac{31}{3}\right)\) 和 \(\left(1,-\frac43\right)\). 因而所求公共点为 \((0,-3)\)，\((-3,-12)\)，\(\left(-2,-\frac{31}{3}\right)\)，\(\left(1,-\frac43\right)\).

\item 令
\[
q(x)=f'(x)=-x^{2}+2bx+c.
\]
由 \(M\) 的定义，至少有
\[
M\geqslant |q(0)|=|c|.
\]
另一方面
\[
q(1)+q(-1)=2(c-1),
\]
故
\[
M\geqslant\max\{|q(1)|,|q(-1)|\}\geqslant |c-1|.
\]
于是
\[
M\geqslant\max\{|c|,|c-1|\}\geqslant\frac12.
\]
取 \(b=0\)，\(c=\frac12\) 时，\(q(x)=\frac12-x^{2}\)，且当 \(-1\leqslant x\leqslant1\) 时，有 \(-\frac12\leqslant q(x)\leqslant\frac12\).
从而 \(M=\frac12\). 因此对任意 \(b\)，\(c\) 恒成立的最大 \(k\) 为
\[
k=\frac12.
\]
\end{enumerate}
\end{solution}
